\subsection{Sums and Direct Sums}
% **Role:** algebraic language of linear structure.
%
% Sums, direct sums, complements, quotient spaces, induced maps, dimension formulas.

Direct sums split spaces into complementary pieces; quotients collapse selected directions. Dimension formulas keep both constructions controlled.

\begin{definition} \hlt{Sums of Subspaces}\\
Let $\mathcal{U}_1,\ldots,\mathcal{U}_m\leq\mathcal{V}$.
\begin{enumerate}[label=(\roman*)]
\setlength{\itemsep}{0pt}
\item The \hlt{sum} of $\mathcal{U}_1,\ldots,\mathcal{U}_m$ is
$\mathcal{U}_1+\cdots+\mathcal{U}_m\defeq\{\mathbf{u}_1+\cdots+\mathbf{u}_m:\mathbf{u}_j\in\mathcal{U}_j\}$.
\item For subspaces $\mathcal{U},\mathcal{W}\leq\mathcal{V}$,
$\mathcal{U}+\mathcal{W}=\operatorname{span}(\mathcal{U}\cup\mathcal{W})=\{\mathbf{u}+\mathbf{w}:\mathbf{u}\in\mathcal{U},\ \mathbf{w}\in\mathcal{W}\}$.
\item The sum is \hlt{direct} if every vector in it has a unique decomposition as $\mathbf{u}_1+\cdots+\mathbf{u}_m$ with $\mathbf{u}_j\in\mathcal{U}_j$; in this case it is denoted $\mathcal{U}_1\oplus\cdots\oplus\mathcal{U}_m$.
\end{enumerate}
\end{definition}

\begin{proposition} \label{prop:sum-subspaces-smallest-containing-subspace} \hlt{Sum of Subspaces Is the Smallest Containing Subspace}\\
Suppose $\mathcal{U}_1,\ldots,\mathcal{U}_m$ are subspaces of $\mathcal{V}$. Then $\mathcal{U}_1+\cdots+\mathcal{U}_m$ is the smallest subspace of $\mathcal{V}$ containing $\mathcal{U}_1,\ldots,\mathcal{U}_m$.
\begin{proof}
Since $\mathbf{0}\in\mathcal{U}_j$ for each $j$, we have $\mathbf{0}\in\mathcal{U}_1+\cdots+\mathcal{U}_m$. If $\mathbf{x}=\mathbf{u}_1+\cdots+\mathbf{u}_m$ and $\mathbf{y}=\mathbf{v}_1+\cdots+\mathbf{v}_m$ lie in $\mathcal{U}_1+\cdots+\mathcal{U}_m$, then for any $c \in \mathbb{F}$, $c\mathbf{x}+\mathbf{y} = (c\mathbf{u}_1+\mathbf{v}_1)+\cdots+(c\mathbf{u}_m+\mathbf{v}_m) \in \mathcal{U}_1+\cdots+\mathcal{U}_m$, because $c\mathbf{u}_j+\mathbf{v}_j\in\mathcal{U}_j$ for each $j$. By Theorem \ref{thm:subspacetest}, $\mathcal{U}_1+\cdots+\mathcal{U}_m$ is a subspace of $\mathcal{V}$.\\
Each $\mathcal{U}_j$ is contained in $\mathcal{U}_1+\cdots+\mathcal{U}_m$ by taking all other summands to be $\mathbf{0}$. If $\mathcal{W}$ is any subspace of $\mathcal{V}$ containing $\mathcal{U}_1,\ldots,\mathcal{U}_m$, then closure under addition gives $\mathbf{u}_1+\cdots+\mathbf{u}_m\in\mathcal{W}$ whenever $\mathbf{u}_j\in\mathcal{U}_j$.\\
Hence $\mathcal{U}_1+\cdots+\mathcal{U}_m\subseteq\mathcal{W}$, so the sum is the smallest such subspace.\\
\hlt{Structural Commentary}: Sums of subspaces are obtained by closing several subspaces under finite addition. Minimality follows because every containing subspace is already closed under those sums.\\
\hlt{Key Insight}: The sum adds no extra structure beyond the linear combinations forced by the given subspaces.
\end{proof}
\end{proposition}

\begin{perspective} \label{per:sum-not-union} \hlt{Sum Is Not Union}\\
The union $\mathcal{U}\cup\mathcal{W}$ keeps vectors from either subspace; the sum $\mathcal{U}+\mathcal{W}$ adds vectors from both. Unless one subspace contains the other, $\mathcal{U}\cup\mathcal{W}$ is not closed under addition.
\end{perspective}

\begin{theorem} \label{thm:finite-union-proper-subspaces-infinite-field} \hlt{Finite Unions of Proper Subspaces}\\
Let $\mathcal{V}\neq\{\mathbf{0}\}$ be over an infinite field $\mathbb{F}$. If $\mathcal{U}_1,\ldots,\mathcal{U}_n\subsetneq\mathcal{V}$ are subspaces, then $\mathcal{V}\neq\bigcup\limits_{j=1}^{n}\mathcal{U}_j$.
\begin{proof}
Assume for contradiction that $\mathcal{V}=\bigcup\limits_{j=1}^{n}\mathcal{U}_j$. After deleting redundant subspaces and reindexing, assume the cover is minimal, so $\mathcal{U}_i\not\subseteq\bigcup\limits_{j\neq i}\mathcal{U}_j$. If $n=1$, then $\mathcal{V}=\mathcal{U}_1$, contradicting $\mathcal{U}_1\subsetneq\mathcal{V}$; hence $n\geq2$.\\
Choose $\mathbf{w}\in\mathcal{U}_1\setminus\bigcup\limits_{j=2}^{n}\mathcal{U}_j$ and $\mathbf{v}\in\mathcal{V}\setminus\mathcal{U}_1$. Since $\mathbf{0}\in\mathcal{U}_j$ for every $j$, we have $\mathbf{w}\neq\mathbf{0}$.\\
Let $A=\{\mathbf{v}+r\mathbf{w}:r\in\mathbb{F}\}$. Then $\abs{A}=\abs{\mathbb{F}}=\infty$, $A\cap\mathcal{U}_1=\emptyset$, and for $j\geq2$, $\abs{A\cap\mathcal{U}_j}\leq1$, since two points $\mathbf{v}+r_1\mathbf{w},\mathbf{v}+r_2\mathbf{w}\in\mathcal{U}_j$ with $r_1\neq r_2$ imply $\mathbf{w}\in\mathcal{U}_j$.\\
Hence $\abs{A\cap\bigcup\limits_{j=1}^{n}\mathcal{U}_j}\leq n-1<\infty$, contradicting $A\subseteq\mathcal{V}=\bigcup\limits_{j=1}^{n}\mathcal{U}_j$.\\
\hlt{Structural Commentary}: Over an infinite field, finite proper subspaces cannot cover a well-chosen affine line.\\
\hlt{Key Insight}: A single one-parameter translate has too many points for finitely many proper intersections.
\end{proof}
\end{theorem}

\begin{definition} \label{def:subspace-meet-and-join} \hlt{Subspace Poset, Meet, and Join}
\begin{enumerate}[label=(\roman*)]
\setlength{\itemsep}{0pt}
\item The \hlt{subspace poset} of $\mathcal{V}$ is
$\operatorname{Sub}(\mathcal{V})\defeq\{\mathcal{U}\subseteq\mathcal{V}:\mathcal{U}\leq\mathcal{V}\}$, ordered by inclusion.
\item For $\{\mathcal{U}_i:i\in I\}\subseteq\operatorname{Sub}(\mathcal{V})$, the \hlt{meet} is
$\bigwedge\limits_{i\in I}\mathcal{U}_i\defeq\bigcap\limits_{i\in I}\mathcal{U}_i$.
\item For $\{\mathcal{U}_i:i\in I\}\subseteq\operatorname{Sub}(\mathcal{V})$, the \hlt{join} is
$\bigvee\limits_{i\in I}\mathcal{U}_i\defeq\sum\limits_{i\in I}\mathcal{U}_i\defeq\operatorname{span}\left(\bigcup\limits_{i\in I}\mathcal{U}_i\right)$.
\item The empty meet is $\mathcal{V}$, and the empty join is $\{\mathbf{0}\}$.
\end{enumerate}
\end{definition}

\begin{proposition} \label{prop:subspaces-form-complete-lattice} \hlt{Subspace Meet and Join}\\
For subspaces $\{\mathcal{U}_i:i\in I\}$ of $\mathcal{V}$, ordered by inclusion, $\bigwedge\limits_{i\in I}\mathcal{U}_i=\inf\{\mathcal{U}_i:i\in I\}$; and $\bigvee\limits_{i\in I}\mathcal{U}_i=\sup\{\mathcal{U}_i:i\in I\}$.
\begin{proof}
If $I=\emptyset$, conventions give vacuous bounds $\mathcal{V}$ and $\{\mathbf{0}\}$. Assume $I\neq\emptyset$.\\
The meet $M=\bigcap_{i\in I}\mathcal{U}_i$ is a subspace by Proposition \ref{prop:intersectionsubspaces}; moreover $M\subseteq\mathcal{U}_i$ for all $i$, and any subspace $\mathcal{W}$ with $\mathcal{W}\subseteq\mathcal{U}_i$ for all $i$ satisfies $\mathcal{W}\subseteq M$. The join $J=\operatorname{span}\left(\bigcup_{i\in I}\mathcal{U}_i\right)$ is a subspace by Theorem \ref{thm:span-smallest-containing-subspace}; moreover $\mathcal{U}_i\subseteq J$ for all $i$, and any subspace $\mathcal{W}$ with $\mathcal{U}_i\subseteq\mathcal{W}$ for all $i$ contains $\bigcup_{i\in I}\mathcal{U}_i$, hence contains $J$.\\
\hlt{Structural Commentary}: Meet enforces common constraints; join closes the combined directions under span.\\
\hlt{Key Insight}: Meet is intersection, while join is span of the union.
\end{proof}
\end{proposition}

\begin{theorem} \label{thm:subspace-lattice-complete-modular} \hlt{Subspace Lattice}\\
The subspace poset $\operatorname{Sub}(\mathcal{V})$ is a complete modular lattice.\\
Explicitly, for every $\{\mathcal{U}_i:i\in I\}\subseteq\operatorname{Sub}(\mathcal{V})$,
$\bot=\{\mathbf{0}\}$, $\top=\mathcal{V}$, $\bigwedge\limits_{i\in I}\mathcal{U}_i=\bigcap\limits_{i\in I}\mathcal{U}_i$, $\bigvee\limits_{i\in I}\mathcal{U}_i=\sum\limits_{i\in I}\mathcal{U}_i=\operatorname{span}\left(\bigcup\limits_{i\in I}\mathcal{U}_i\right)$.\\
Moreover, if $\mathcal{U}\subseteq\mathcal{S}$, then $\mathcal{S}\cap(\mathcal{T}+\mathcal{U})=(\mathcal{S}\cap\mathcal{T})+\mathcal{U}$.
\begin{proof}
Proposition \ref{prop:subspaces-form-complete-lattice} gives arbitrary meets and joins with the displayed formulas; the empty meet and join give $\top=\mathcal{V}$ and $\bot=\{\mathbf{0}\}$. It remains to prove modularity. Assume $\mathcal{U}\subseteq\mathcal{S}$. If $\mathbf{x}\in(\mathcal{S}\cap\mathcal{T})+\mathcal{U}$, then $\mathbf{x}=\mathbf{y}+\mathbf{u}$ with $\mathbf{y}\in\mathcal{S}\cap\mathcal{T}$ and $\mathbf{u}\in\mathcal{U}$, hence $\mathbf{x}\in\mathcal{S}\cap(\mathcal{T}+\mathcal{U})$. Conversely, if $\mathbf{x}\in\mathcal{S}\cap(\mathcal{T}+\mathcal{U})$, write $\mathbf{x}=\mathbf{t}+\mathbf{u}$ with $\mathbf{t}\in\mathcal{T}$ and $\mathbf{u}\in\mathcal{U}$. Since $\mathbf{x}\in\mathcal{S}$ and $\mathbf{u}\in\mathcal{U}\subseteq\mathcal{S}$, $\mathbf{t}=\mathbf{x}-\mathbf{u}\in\mathcal{S}\cap\mathcal{T}$, so $\mathbf{x}\in(\mathcal{S}\cap\mathcal{T})+\mathcal{U}$.\\
\hlt{Structural Commentary}: The subspace lattice is complete and modular, not generally distributive.\\
\hlt{Key Insight}: If one summand already lies inside $\mathcal{S}$, intersecting with $\mathcal{S}$ can be pushed through the sum.
\end{proof}
\end{theorem}

\begin{figure}[H]
  \centering
  \begin{tikzpicture}[
  x=0.88cm,
  y=0.92cm,
  >=Stealth,
  every node/.style={font=\small},
  point/.style={circle, fill=gray!75, inner sep=1.55pt},
  edge/.style={->, gray!65!black, line width=0.7pt},
  hint/.style={font=\scriptsize, text=gray!70!black, align=center, inner sep=1pt}
]
  \node[point] (bot) at (-3.15,0) {};
  \node[point] (meet) at (-1.55,0) {};
  \node[point] (u) at (0,0.42) {};
  \node[point] (w) at (0,-0.42) {};
  \node[point] (join) at (1.55,0) {};
  \node[point] (top) at (3.15,0) {};

  \draw[edge] (bot) -- (meet);
  \draw[edge] (meet) -- (u);
  \draw[edge] (meet) -- (w);
  \draw[edge] (u) -- (join);
  \draw[edge] (w) -- (join);
  \draw[edge] (join) -- (top);

  \node at (-3.15,-0.36) {$\{\mathbf{0}\}$};
  \node at (-1.55,-0.36) {$\mathcal{U}\cap\mathcal{W}$};
  \node[hint] at (-1.55,-0.68) {meet};
  \node at (0,0.78) {$\mathcal{U}$};
  \node at (0,-0.78) {$\mathcal{W}$};
  \node at (1.55,0.36) {$\mathcal{U}+\mathcal{W}$};
  \node[hint] at (1.55,-0.36) {join};
  \node at (3.15,-0.36) {$\mathcal{V}$};
\end{tikzpicture}

  \caption{For two subspaces, intersection is the meet and sum is the join under inclusion.}
  \label{fig:subspace-lattice-meet-join}
\end{figure}

\begin{perspective} \label{per:sum-versus-direct-sum-existence} \hlt{Sums Always Exist; Directness Is Extra}\\
For subspaces $\mathcal{U},\mathcal{W}\subseteq\mathcal{V}$, the sum $\mathcal{U}+\mathcal{W}$ is always a subspace by Proposition \ref{prop:sum-subspaces-smallest-containing-subspace}. Writing $\mathcal{U}\oplus\mathcal{W}$ asserts the extra condition $\mathcal{U}\cap\mathcal{W}=\{\mathbf{0}\}$, equivalently uniqueness of the decomposition $\mathbf{u}+\mathbf{w}$.
\end{perspective}

\begin{proposition} \label{prop:uniqueness-direct-sum-two-subspaces} \hlt{Two-Subspace Directness Criterion}\\
Let $\mathcal{U}$ and $\mathcal{W}$ be subspaces of $\mathcal{V}$. Then $\mathcal{U}+\mathcal{W}$ is direct if and only if $\mathcal{U}\cap\mathcal{W}=\{\mathbf{0}\}$.
\begin{proof}
\begin{enumerate}[label=(\roman*)]
\setlength{\itemsep}{0pt}
\item[$(\Rightarrow)$] Suppose $\mathcal{U}+\mathcal{W}$ is direct. If $\mathbf{x}\in\mathcal{U}\cap\mathcal{W}$, then $\mathbf{x}=\mathbf{x}+\mathbf{0}=\mathbf{0}+\mathbf{x}$ gives two decompositions of $\mathbf{x}$ with one term in each subspace. Uniqueness gives $\mathbf{x}=\mathbf{0}$, so $\mathcal{U}\cap\mathcal{W}=\{\mathbf{0}\}$.
\item[$(\Leftarrow)$] Suppose $\mathcal{U}\cap\mathcal{W}=\{\mathbf{0}\}$. Every vector in $\mathcal{U}+\mathcal{W}$ has at least one representation $\mathbf{u}+\mathbf{w}$. If $\mathbf{u}_1+\mathbf{w}_1=\mathbf{u}_2+\mathbf{w}_2$, with $\mathbf{u}_1,\mathbf{u}_2\in\mathcal{U}$ and $\mathbf{w}_1,\mathbf{w}_2\in\mathcal{W}$, then $\mathbf{u}_1-\mathbf{u}_2=\mathbf{w}_2-\mathbf{w}_1\in\mathcal{U}\cap\mathcal{W}$, hence equals $\mathbf{0}$.\\
Thus $\mathbf{u}_1=\mathbf{u}_2$ and $\mathbf{w}_1=\mathbf{w}_2$, so the sum is direct.
\end{enumerate}
\hlt{Structural Commentary}: Trivial intersection prevents cancellation between summands.\\
\hlt{Key Insight}: Directness means unique decomposition.
\end{proof}
\end{proposition}

\begin{proposition} \label{prop:direct-sum-equivalent-conditions} \hlt{Equivalent Conditions for a Direct Sum}\\
For subspaces $\mathcal{U}_1, \ldots, \mathcal{U}_m \subseteq \mathcal{V}$, the following are equivalent:
\begin{enumerate}[label=(\roman*)]
\setlength{\itemsep}{0pt}
\item $\mathcal{V} = \mathcal{U}_1 \oplus \cdots \oplus \mathcal{U}_m$;
\item every $\mathbf{v} \in \mathcal{V}$ has a unique representation $\mathbf{v} = \mathbf{u}_1 + \cdots + \mathbf{u}_m$ with $\mathbf{u}_j \in \mathcal{U}_j$;
\item $\mathcal{V}=\mathcal{U}_1+\cdots+\mathcal{U}_m$, and if $\mathbf{u}_1 + \cdots + \mathbf{u}_m = \mathbf{0}$ with $\mathbf{u}_j \in \mathcal{U}_j$, then $\mathbf{u}_1 = \cdots = \mathbf{u}_m = \mathbf{0}$.
\end{enumerate}
\begin{proof}
\begin{enumerate}[label=(\roman*)]
\setlength{\itemsep}{0pt}
\item[$(i)$]$\Rightarrow (ii)$ This follows from the definition of direct sum.
\item[$(ii)$]$\Rightarrow (iii)$ Every vector lies in $\mathcal{U}_1+\cdots+\mathcal{U}_m$, so $\mathcal{V}=\mathcal{U}_1+\cdots+\mathcal{U}_m$. If $\mathbf{u}_1+\cdots+\mathbf{u}_m=\mathbf{0}$, then this representation of $\mathbf{0}$ must equal $\mathbf{0}+\cdots+\mathbf{0}$ by uniqueness, so each $\mathbf{u}_j=\mathbf{0}$.
\item[$(iii)$]$\Rightarrow (i)$ Since $\mathcal{V}=\mathcal{U}_1+\cdots+\mathcal{U}_m$, every $\mathbf{v}\in\mathcal{V}$ has at least one representation.\\
If $\mathbf{v}=\mathbf{u}_1+\cdots+\mathbf{u}_m=\mathbf{w}_1+\cdots+\mathbf{w}_m$, then $(\mathbf{u}_1-\mathbf{w}_1)+\cdots+(\mathbf{u}_m-\mathbf{w}_m)=\mathbf{0}$. By (iii), $\mathbf{u}_j-\mathbf{w}_j=\mathbf{0}$ for every $j$, so the representation is unique. Hence $\mathcal{V}=\mathcal{U}_1\oplus\cdots\oplus\mathcal{U}_m$.
\end{enumerate}
\hlt{Structural Commentary}: Direct sums require both spanning by the summands and no cancellation among them. The zero-representation condition isolates the no-cancellation part.\\
\hlt{Key Insight}: Uniqueness is equivalent to saying the only way the summands add to zero is trivially.
\end{proof}
\end{proposition}

\begin{theorem} \label{thm:dimension-of-product} \hlt{Dimension of a Product}\\
If $\mathcal{V}_1,\ldots,\mathcal{V}_m$ are finite-dimensional vector spaces over $\mathbb{F}$, then $\mathcal{V}_1\times\cdots\times\mathcal{V}_m$ is finite-dimensional,\\
and $ \dim(\mathcal{V}_1\times\cdots\times\mathcal{V}_m)=\dim(\mathcal{V}_1)+\cdots+\dim(\mathcal{V}_m)$.
\begin{proof}
For each $j$, let $(\mathbf{v}_{j,1},\ldots,\mathbf{v}_{j,n_j})$ be a basis of $\mathcal{V}_j$. For every $j$ and $k$, let $\mathbf{e}_{j,k}$ denote the element of $\mathcal{V}_1\times\cdots\times\mathcal{V}_m$ whose $j$th component is $\mathbf{v}_{j,k}$ and whose other components are zero.\\
The list of all $\mathbf{e}_{j,k}$ spans the product because each component can be expanded in its own basis. If a linear combination of the $\mathbf{e}_{j,k}$ is zero, then looking at the $j$th component gives a linear relation among $(\mathbf{v}_{j,1},\ldots,\mathbf{v}_{j,n_j})$, so all coefficients in that component are zero. Hence the list is linearly independent and therefore is a basis of the product. Its length is $n_1+\cdots+n_m$.\\
\hlt{Structural Commentary}: A basis for a product is formed by placing each component basis into its own slot.\\
\hlt{Key Insight}: Product dimension counts the independent directions contributed by each factor.
\end{proof}
\end{theorem}

\begin{definition} \label{def:external-product-family} \hlt{External Product}\\
Let $\{\mathcal{V}_i:i\in I\}$ be $\mathbb{F}$-vector spaces.
\begin{enumerate}[label=(\roman*)]
\setlength{\itemsep}{0pt}
\item The \hlt{external direct product} is $\prod\limits_{i\in I}\mathcal{V}_i =\{\mathbf{v}:I\to\bigcup\limits_{i\in I}\mathcal{V}_i:\mathbf{v}(i)\in\mathcal{V}_i\}$.
\item Operations are componentwise: $(\mathbf{u}+\mathbf{v})(i)=\mathbf{u}(i)+\mathbf{v}(i)$ and $(a\mathbf{v})(i)=a\mathbf{v}(i)$.
\end{enumerate}
\end{definition}

\begin{definition} \label{def:support-and-finite-family-sums} \hlt{Support and Finite Sums of a Family}\\
Let $I$ be an index set. For each $i\in I$, let $\mathcal{X}_i$ be a vector space with zero vector $\mathbf{0}_i$.
\begin{enumerate}[label=(\roman*)]
\setlength{\itemsep}{0pt}
\item \hlt{Support.} For $\mathbf{x}=(\mathbf{x}_i)_{i\in I}\in\prod\limits_{i\in I}\mathcal{X}_i$, define $\operatorname{supp}(\mathbf{x})\defeq\{i\in I:\mathbf{x}_i\neq\mathbf{0}_i\}$.\\
\item The indexed family $\mathbf{x}$ is \hlt{finitely supported} iff $\abs{\operatorname{supp}(\mathbf{x})}<\infty$.
\item \hlt{Finite sum.} If all $\mathcal{X}_i$ lie in a common vector space and $\mathbf{x}$ is finitely supported, define $\sum\limits_{i\in I}\mathbf{x}_i\defeq\sum\limits_{i\in\operatorname{supp}(\mathbf{x})}\mathbf{x}_i$.
\end{enumerate}
\end{definition}

\begin{definition} \label{def:external-products-and-direct-sums-family} \hlt{External Direct Sum}\\
Let $\{\mathcal{V}_i:i\in I\}$ be $\mathbb{F}$-vector spaces. The \hlt{external direct sum} is $\bigoplus\limits_{i\in I}^{\mathrm{ext}}\mathcal{V}_i =\{\mathbf{v}\in\prod\limits_{i\in I}\mathcal{V}_i:\abs{\operatorname{supp}(\mathbf{v})}<\infty\}$. If $\abs{I}<\infty$, then $\bigoplus\limits_{i\in I}^{\mathrm{ext}}\mathcal{V}_i=\prod\limits_{i\in I}\mathcal{V}_i$.
\end{definition}

\begin{proposition} \label{prop:external-direct-sum-family-vector-space} \hlt{External Sums Are Finite-Support Products}\\
With componentwise operations, $\prod\limits_{i\in I}\mathcal{V}_i$ is an $\mathbb{F}$-vector space, $\bigoplus\limits_{i\in I}^{\mathrm{ext}}\mathcal{V}_i =\{\mathbf{v}\in\prod\limits_{i\in I}\mathcal{V}_i:\abs{\operatorname{supp}(\mathbf{v})}<\infty\} \leq \prod\limits_{i\in I}\mathcal{V}_i$.
\begin{proof}
Let $P=\prod_{i\in I}\mathcal{V}_i$. For $\mathbf{u},\mathbf{v}\in P$ and $a,b\in\mathbb{F}$, $(a\mathbf{u}+b\mathbf{v})(i)=a\mathbf{u}(i)+b\mathbf{v}(i)\in\mathcal{V}_i$.\\
The zero $0_P(i)=0_i$, inverses $(-\mathbf{v})(i)=-\mathbf{v}(i)$, and all axioms are coordinatewise.\\
Let $S=\bigoplus_{i\in I}^{\mathrm{ext}}\mathcal{V}_i$. If $\mathbf{u},\mathbf{v}\in S$ and $a,b\in\mathbb{F}$, then $\operatorname{supp}(a\mathbf{u}+b\mathbf{v})\subseteq \operatorname{supp}(\mathbf{u})\cup\operatorname{supp}(\mathbf{v})$, so $a\mathbf{u}+b\mathbf{v}\in S$.\\
Also $0_P\in S$. Hence $S\leq P$.\\
\hlt{Structural Commentary}: $\prod_i\mathcal{V}_i$: arbitrary support; $\bigoplus_i^{\mathrm{ext}}\mathcal{V}_i$: finite support.\\
\hlt{Key Insight}: Finite support is exactly what finite linear combinations preserve.
\end{proof}
\end{proposition}

\begin{definition} \label{def:internal-direct-sum-family} \hlt{Internal Direct Sum of a Family}\\
For $\{\mathcal{U}_i:i\in I\}\subseteq\operatorname{Sub}(\mathcal{V})$, $\mathcal{V}=\bigoplus\limits_{i\in I}\mathcal{U}_i$ means every $\mathbf{v}\in\mathcal{V}$ has a unique finitely supported $\mathbf{u}=(\mathbf{u}_i)_{i\in I}$ with $\mathbf{u}_i\in\mathcal{U}_i$ and $\mathbf{v}=\sum\limits_{i\in I}\mathbf{u}_i$.\\
Equivalently, $\mathcal{V}=\sum\limits_{i\in I}\mathcal{U}_i$ and every finite relation $\sum\limits_{i\in F}\mathbf{u}_i=\mathbf{0}$, with $F\subseteq I$ finite and $\mathbf{u}_i\in\mathcal{U}_i$, has $\mathbf{u}_i=\mathbf{0}$ $\forall i\in F$.
\end{definition}

\begin{proposition} \label{prop:internal-direct-sum-family-criterion} \hlt{Family Direct-Sum Criterion}\\
For $\{\mathcal{U}_i\}_{i\in I}\subseteq\operatorname{Sub}(\mathcal{V})$, $\sum\limits_{i\in I}\mathcal{U}_i$ is direct $\Longleftrightarrow$ $\forall i\in I$, $\mathcal{U}_i\cap\sum\limits_{j\in I\setminus\{i\}}\mathcal{U}_j=\{\mathbf{0}\}$.
\begin{proof}
\begin{enumerate}[label=(\roman*)]
\setlength{\itemsep}{0pt}
\item[$(\Rightarrow)$] If $\mathbf{x}\in\mathcal{U}_i\cap\sum\limits_{j\in I\setminus\{i\}}\mathcal{U}_j$, then $\mathbf{x}=\sum\limits_{j\in F}\mathbf{u}_j$ for some finite $F\subseteq I\setminus\{i\}$ and $\mathbf{u}_j\in\mathcal{U}_j$.\\
Directness applied to $\mathbf{x}-\sum\limits_{j\in F}\mathbf{u}_j=\mathbf{0}$ gives $\mathbf{x}=\mathbf{0}$.
\item[$(\Leftarrow)$] If $\sum\limits_{k=1}^{m}\mathbf{u}_{i_k}=\mathbf{0}$ with distinct $i_k$ and $\mathbf{u}_{i_k}\in\mathcal{U}_{i_k}$, then for each $k$, $\mathbf{u}_{i_k}=-\sum\limits_{\ell\neq k}\mathbf{u}_{i_\ell}\in\mathcal{U}_{i_k}\cap\sum\limits_{j\in I\setminus\{i_k\}}\mathcal{U}_j=\{\mathbf{0}\}$.\\
Thus all finite zero-relations are trivial, so the sum is direct.
\end{enumerate}
\hlt{Structural Commentary}: Family directness requires each summand to miss all others together.\\
\hlt{Key Insight}: Directness is subspace independence.
\end{proof}
\end{proposition}


\begin{proposition} \label{prop:direct-sum-map-injective} \hlt{Direct Sum Map}\\
Suppose $\mathcal{U}_1,\ldots,\mathcal{U}_m$ are subspaces of $\mathcal{V}$. Define $\Gamma:\mathcal{U}_1\times\cdots\times\mathcal{U}_m\to\mathcal{U}_1+\cdots+\mathcal{U}_m$ by $\Gamma(\mathbf{u}_1,\ldots,\mathbf{u}_m)=\mathbf{u}_1+\cdots+\mathbf{u}_m$.\\
Then $\mathcal{U}_1+\cdots+\mathcal{U}_m$ is a direct sum if and only if $\Gamma$ is injective.
\begin{proof}
\begin{enumerate}[label=(\roman*)]
\setlength{\itemsep}{0pt}
\item[$(\Rightarrow)$] Suppose $\mathcal{U}_1+\cdots+\mathcal{U}_m$ is direct. If $\Gamma(\mathbf{u}_1,\ldots,\mathbf{u}_m)=\mathbf{0}$, then $\mathbf{u}_1+\cdots+\mathbf{u}_m=\mathbf{0}$, so directness gives $\mathbf{u}_j=\mathbf{0}$ for all $j$. Hence $\ker\Gamma=\{(\mathbf{0},\ldots,\mathbf{0})\}$, and $\Gamma$ is injective.
\item[$(\Leftarrow)$] Suppose $\Gamma$ is injective. If $\mathbf{u}_1+\cdots+\mathbf{u}_m=\mathbf{0}$ with $\mathbf{u}_j\in\mathcal{U}_j$, then $\Gamma(\mathbf{u}_1,\ldots,\mathbf{u}_m)=\Gamma(\mathbf{0},\ldots,\mathbf{0})$. Injectivity gives $\mathbf{u}_j=\mathbf{0}$ for all $j$, so every zero decomposition is trivial. Hence the sum is direct.
\end{enumerate}
\hlt{Structural Commentary}: The addition map identifies decompositions precisely through its kernel.\\
\hlt{Key Insight}: Directness is injectivity of $\Gamma$.
\end{proof}
\end{proposition}

\begin{figure}[H]
  \centering
  \begin{tikzpicture}[
  >=Stealth,
  every node/.style={font=\small},
  ambient/.style={draw=gray!45, fill=gray!3, rounded corners=3pt, line width=0.55pt,
    minimum width=4.55cm, minimum height=2.5cm},
  set/.style={draw=gray!55, fill=gray!6, line width=0.55pt},
  box/.style={draw=gray!45, fill=gray!3, rounded corners=3pt, line width=0.55pt,
    minimum width=2.45cm, minimum height=0.78cm, align=center},
  map/.style={->, gray!65!black, line width=0.75pt}
]
  \node[ambient] at (-2.15,0) {};
  \node[anchor=north west] at (-4.18,1.08) {$\mathcal{U}+\mathcal{W}$};

  \draw[set, fill=blue!7] (-2.75,0) ellipse (1.05cm and 0.65cm);
  \draw[set, fill=teal!8, fill opacity=0.75] (-1.55,0) ellipse (1.05cm and 0.65cm);
  \node at (-3.15,0.38) {$\mathcal{U}$};
  \node at (-1.14,0.38) {$\mathcal{W}$};
  \node at (-2.15,-0.02) {$\mathcal{U}\cap\mathcal{W}$};

  \node[box] (prod) at (2.25,0.55) {$\mathcal{U}\times\mathcal{W}$};
  \node[box] (sum) at (2.25,-0.75) {$\mathcal{U}+\mathcal{W}$};
  \draw[map] (prod) -- (sum);
  \node[gray!70!black, anchor=west] at (2.68,-0.10) {$\Gamma(\mathbf{u},\mathbf{w})=\mathbf{u}+\mathbf{w}$};
\end{tikzpicture}

  \caption{Intersection, sum, and product of subspaces.}
  \label{fig:subspace-sum-product-schematic}
\end{figure}
